\documentclass[10pt,a4paper]{amsart}
\usepackage[margin=19mm]{geometry}
\usepackage{amsmath,amssymb,mathtools}
\usepackage[scr=rsfs,cal=euler]{mathalfa}
\usepackage{xfrac}
\usepackage[unicode,hidelinks,bookmarks=false]{hyperref}
\newcommand{\ie}{i.e.\ }
\newcommand{\cf}{cf.\ }
\newcommand{\resp}{respectively\ }
\newcommand{\Subsection}{\subsection}
\providecommand{\nobreakdash}{\nobreak\hbox{-}}
\setlength{\emergencystretch}{2em}
\multlinegap0pt
\usepackage{bm}
\usepackage{bbm}
\usepackage{dsfont}
\usepackage{xspace}
\usepackage{cancel}
\usepackage{mathtools}
\usepackage[shortlabels]{enumitem}
\usepackage{amsthm}
\makeatletter
\@ifundefined{theorem}{\newtheorem{theorem}{Theorem}}{}
\makeatother
\newcommand{\cM}{\mathcal{M}}
\newcommand{\funF}{\mathfrak{F}}
\renewcommand{\hat}{\widehat}
\title[Generalized rational Prony and Bernoulli methods]{Generalized rational Prony and Bernoulli methods}
\author{Tam\'as D\'ozsa, Matthias Voigt, Zolt\'an Szab\'o, J\'ozsef Bokor, P\'eter Kov\'acs}
\date{Source: arXiv:2510.03510v1 (CC BY 4.0)}
\begin{document}
\maketitle

\noindent\emph{Source and licence.} Generalized rational Prony and Bernoulli methods. Version arXiv:2510.03510v1 (CC BY 4.0), distributed under the Creative Commons Attribution 4.0 International licence, as recorded in the arXiv licence metadata for this exact version and shown on the abstract page https://arxiv.org/abs/2510.03510v1. Full rights notice: licenses/arxiv-2510.03510-CC-BY-4.0.txt.\par\smallskip

\noindent\textit{Editorial scope.} Counted unit: the Theorem whose source label is \texttt{thm:ratpron} in the article (source0.tex lines 505--526, source label \texttt{thm:ratpron}), delivered together with the article's own complete original proof of it (source0.tex lines 528--570, one \texttt{proof} environment of 43 lines). No mathematics is written, repaired or re-derived: statement and proof are the article's own text line for line. Required context retained uncounted: eq:ma (lines 129--132); eq:def\_rlambda (lines 299--301). Every definition, display and label of those lines is retained unchanged. Omitted from this package and not redistributed: 1--128, 133--298, 302--504, 527--527, 571--986. Nothing inside a retained excerpt is omitted or altered: every retained body is delivered exactly as the article's lines are written, so no omission receipt of any kind accompanies this package. Citations: the retained excerpts carry no citation command at all, so no bibliography entry of the article is transcribed and no reference is redistributed. Reference adaptation: none was needed and none was made; every reference inside the delivered unit resolves inside it, so no unresolved reference or citation is produced. Printed slips kept literal: no printed word, symbol, hypothesis or number of the counted statement or of its proof was altered. Layout and class: the article's own class is \texttt{arxiv} with packages including natbib, mathtools, amsmath, amsfonts, amsthm, enumitem, booktabs, pgfplots, tikz; the wrapper is the lane's pinned \texttt{amsart} editorial wrapper at 10pt on A4, which restates the theorem-like environments the retained text needs and adds the article's own macro definitions that the retained text needs (\texttt{\textbackslash cM}, \texttt{\textbackslash funF}, \texttt{\textbackslash hat}), with extra wrapper packages (bm, bbm, dsfont, xspace, cancel, mathtools, enumitem). The delivered numbering is the unit's own. Assets: no retained span contains a figure environment, a graphics-inclusion command, a PDF-page inclusion, a table or a TikZ picture; no image file is copied into this package, the package declares no asset dependency and no image file is redistributed with it. Licence: version arXiv:2510.03510v1 of this article is distributed under the Creative Commons Attribution 4.0 International licence, as recorded in the arXiv licence metadata for this exact version and shown on the abstract page https://arxiv.org/abs/2510.03510v1. A full rights notice is delivered at licenses/arxiv-2510.03510-CC-BY-4.0.txt. Characters: every retained excerpt is pure ASCII, so the delivered engine, XeLaTeX with its default Latin Modern text font, reports no missing character.
\begin{equation}
     \label{eq:ma}
     \cM := \mathcal{M}(A,\Lambda) := \operatorname{span} \{v_{\lambda_k}:\ k=1,\ldots,M \} = \left\{ f = \sum_{k=1}^M c_k v_{\lambda_k}: \ c_k \in \mathbb{C}\right\}.
\end{equation}

\begin{equation}\label{eq:def_rlambda}
r_{\lambda_k}(z) := \frac{1}{1-\overline{\lambda_k}z} \quad (z\in \mathbb{D}, \ k=1,\ldots,M).
\end{equation}

\begin{theorem}[Generalized operator-based Prony methods and pole finding problems]
\label{thm:ratpron}
    Let the signal space $\cM$ be defined as  in~\eqref{eq:ma} and let $f \in \cM$, i.e.,
    \begin{equation*}
        f = \sum_{k=1}^{M} c_k v_{\lambda_k}.
    \end{equation*}
   % and denote the set of active eigenvalues of $A$ by $\Lambda := \{ \lambda_1, \ldots, \lambda_M \} \subset \mathbb{C}$. 
   Then the following statements are satisfied:
    \begin{enumerate}[label = \alph*)]
    \item There exists a linear operator $\mathcal{G} : \cM \to H_2(\mathbb{D})$ such that
    \begin{equation}
        \label{eq:G}
        \mathcal{G}f := \sum_{k=1}^{M} \mu(c_k) r_{\varrho(\lambda_k)}, 
    \end{equation}
    where $\mu : \mathbb{C} \to \mathbb{C}$ and $\varrho: \mathbb{C} \to \mathbb{D}$ are injective functions and $r_{\varrho(\lambda_k)} \in H_2(\mathbb{D})$ with  $r_{\varrho(\lambda_k)}(z) := \frac{1}{1-\overline{\varrho(\lambda_k)}z}$ for all $k=1,\ldots,M$ (cf. Eq.~\eqref{eq:def_rlambda}).
    \item Assume that we have given a canonical evaluation scheme $\mathcal{F} = {(\mathcal{F}_m)}_{m \ge 0} = {(\funF \circ \Psi_\varphi^m)}_{m \ge 0}$ with an iteration operator $\Psi_\varphi:\cM \to\cM$. Then, $\mathcal{G}$ can be expressed explicitly as
    $$
        \mathcal{G}f := \mathcal{Z}_w[\mathcal{F}f],
    $$
    for a sufficiently large $w > 0$ and where $\mathcal{F}f:=(\mathcal{F}_0f,\mathcal{F}_1f,\ldots)$.
    \end{enumerate}
\end{theorem}

\begin{proof}
First we show the existence of such an operator, then we define it in a constructive manner without assuming a priori knowledge of the parameters $\lambda_k$ and $c_k$.
\begin{enumerate}[label = \alph*)]
\item  A linear mapping $\mathcal{G} : \cM \to H_2(\mathbb{D})$ which satisfies Eq.~\eqref{eq:G} exists. Indeed, since $\Lambda$ is bounded, there exists a $C_{\Lambda} > 0$ such that
\begin{equation}
    \label{eq:cf}
    \max_{k=1,\ldots,M}|\lambda_k| < C_{\Lambda}.
\end{equation}
Then, we can define $\mathcal{G} : \cM \to H_2(\mathbb{D})$ as
\begin{equation*}
    (\mathcal{G}f)(z) := \sum_{k=1}^{M} c_k r_{\lambda_k / C_{\Lambda}}(z) = \sum_{k=1}^{M} \frac{c_k}{1 - \overline{\lambda_k/ C_{\Lambda}}z} =: \hat{G}(z).
\end{equation*}
Then, $\hat{G}$ clearly belongs to an $M$-dimensional subspace of $H_2(\mathbb{D})$ and the linearity of $\mathcal{G}$ is obvious. 

\item Next, we define $\mathcal{G}$ in a constructive manner. % Suppose a canonical evaluation scheme in $V$ is available. 
%Recall, that a canonical evaluation scheme takes the form of
%\begin{equation}
%    \mathcal{F}_m := \funF \circ \Psi_{\varphi}^m,
%\end{equation}
%where $\Psi_\varphi : \cM \to \cM$ with an appropriate injective function $\varphi$ is an iteration operator (see Definition~\ref{def:iterop}). 
Consider the sequence $g^f := \big(g_m^f\big)_{m \ge 0}$ with 
\begin{equation*}
    g^{f}_m := \mathcal{F}_m(f) = \sum_{k=1}^M c_k \funF( \Psi_\varphi^m v_{\lambda_k} ) = \sum_{k=1}^M c_k \varphi(\lambda_k)^m \funF( v_{\lambda_k}) = \sum_{k=1}^{M} \mu(c_k) \varphi(\lambda_k)^m\quad (m \in \mathbb{N} \cup \{0\}),
\end{equation*}
where $\mu(c_k) := c_k \funF(v_{\lambda_k})$. 
%We note that (see also the discussion under Definition~\ref{def:evalscheme}) we expect a canonical evaluation scheme to satisfy $\funF (v_{\lambda_k}) \neq 0$ for $k=1,\ldots,M$.

Choosing $w > \max_{k=1,\ldots,M} |\varphi(\lambda_k)|$, for $z \in \overline{\mathbb{D}}$ we get
\begin{align*}
    \label{eq:h2form}
         G(z) &= \mathcal{Z}_{w} \big[g^f\big] (z) = \sum_{n=0}^{\infty} g^f_n \left( \frac{z}{w} \right)^n = \sum_{k=1}^{M} \mu(c_k) \sum_{n=0}^{\infty} \left( \frac{\varphi(\lambda_k)}{w} \right)^n z^n =\sum_{k=1}^{M} \frac{\mu(c_k)}{1 - {\frac{\varphi(\lambda_k)}{w}}z},
\end{align*}
where in the notation of the theorem, we have $\varrho(\lambda) = \overline{\frac{\varphi(\lambda)}{w}}$ for all $\lambda \in \mathbb{C}$.
%Here the transformation $\mathcal{Z}_{w}$ was defined according to Eq.~\eqref{eq:zw}.
The function $G$ clearly belongs to $H_2(\mathbb{D})$. Using the above, the explicit form of the transformation $\mathcal{G}$ can be written as
\begin{equation*}
 \mathcal{G}f = G = \mathcal{Z}_w[\mathcal{F}f].
   % G(z) := \left(\mathcal{Z}_{w} \circ \mathcal{F}\right)f,
\end{equation*}
\end{enumerate}

%$: V \to \ell$ denotes the functional sequence corresponding of the canonical evaluation scheme $\mathcal{F} = {(\mathcal{F}_m)}_{m \ge 0}$.
\end{proof}

\end{document}
